\documentclass[12pt,letterpaper]{article}
% Common article worksheet layout. Load after \documentclass.
\usepackage[margin=0.5in,top=0.85in,headheight=0pt,headsep=0pt,footskip=0pt]{geometry}
\usepackage{amsmath,amssymb}
\usepackage{tikz,graphicx}
\definecolor{mtuGreen}{RGB}{0,79,42}
\definecolor{mtuGreenLight}{RGB}{60,133,90}
\definecolor{mtuGreenPale}{RGB}{223,236,227}
\newcommand{\vect}[1]{\mathbf{#1}}
\setlength{\parindent}{0pt}
\setlength{\parskip}{8pt}
\pagestyle{empty}
% Keep all marks inside a half-inch print margin, including their strokes.
% Odd pages are left-aligned; even pages are right-aligned.
\AddToHook{shipout/foreground}{%
  \begin{tikzpicture}[remember picture,overlay]
    \ifodd\value{page}
      \node[anchor=north west,draw=black,rectangle,line width=1.5pt,inner xsep=3pt,inner ysep=1pt,font=\fontsize{16}{16}\selectfont\bfseries]
        at ([xshift=0.52in,yshift=-0.52in]current page.north west) {\thepage};
    \else
      \node[anchor=north east,draw=black,rectangle,line width=1.5pt,inner xsep=3pt,inner ysep=1pt,font=\fontsize{16}{16}\selectfont\bfseries]
        at ([xshift=-0.52in,yshift=-0.52in]current page.north east) {\thepage};
    \fi
    % First-page marker matches the dimensions of the boxed page number 1.
    \ifnum\value{page}=1\relax
      \node[anchor=north east,draw=black,fill=black,rectangle,line width=1.5pt,inner xsep=3pt,inner ysep=1pt,font=\fontsize{16}{16}\selectfont\bfseries]
        at ([xshift=-0.52in,yshift=-0.52in]current page.north east) {\phantom{1}};
    \fi
  \end{tikzpicture}%
}
% Fixed name-box dimensions and page position shared by every worksheet.
\newcommand{\WorksheetNameBox}{%
  \begingroup\setlength{\fboxsep}{3pt}\setlength{\fboxrule}{0.4pt}%
  \fbox{\rule[-0.30in]{0pt}{0.42in}\makebox[3.3in][l]{Name:}}%
  \endgroup}
\newcommand{\WorksheetHeading}[1]{%
  \noindent{\small\textbf{MA3521 --- #1}}\hfill\WorksheetNameBox\par}
% A blank entry with fixed writing room, usable inside a matrix.
\newcommand{\AnswerEntry}{\rule[-0.17in]{0pt}{0.47in}\makebox[0.67in]{}}
\newcommand{\AnswerVector}{\begin{pmatrix}\AnswerEntry\\\AnswerEntry\\\AnswerEntry\end{pmatrix}}
\newcommand{\AnswerMatrix}{\begin{pmatrix}
  \AnswerEntry&\AnswerEntry&\AnswerEntry\\
  \AnswerEntry&\AnswerEntry&\AnswerEntry\\
  \AnswerEntry&\AnswerEntry&\AnswerEntry
\end{pmatrix}}

% Reusable three-tank network. Load once in the preamble.
% \ThreeTankDiagram{capacity1=100,volume1=60,concentration1=2,...}
% Volumes: L. Concentrations: g/L. Flow rates: L/min.
\usetikzlibrary{arrows.meta}
\pgfkeys{/ode tanks/.is family,/ode tanks,
  capacity1/.store in=\ODEcapone, capacity2/.store in=\ODEcaptwo, capacity3/.store in=\ODEcapthree,
  volume1/.store in=\ODEvolone, volume2/.store in=\ODEvoltwo, volume3/.store in=\ODEvolthree,
  concentration1/.store in=\ODEcone, concentration2/.store in=\ODEctwo, concentration3/.store in=\ODEcthree,
  inflow/.store in=\ODEqin, inflowconcentration/.store in=\ODEcin,
  flow12/.store in=\ODEqonetwo, flow21/.store in=\ODEqtwoone,
  flow23/.store in=\ODEqtwothree, flow32/.store in=\ODEqthreetwo,
  outflow/.store in=\ODEqout,
  saltRate1/.store in=\ODEsaltone, saltRate2/.store in=\ODEsalttwo,
  saltRate3/.store in=\ODEsaltthree}

% Shared compact labels used by both lectures and worksheets.
\newcommand{\ODETankHeading}[2]{%
  \node[font=\bfseries] at (0,3.12) {$T_{#1}$ (#2 L)};}
\newcommand{\ODETankLabels}[4]{%
  \pgfmathtruncatemacro{\InitialSalt}{#3*#4}
  \node[anchor=north,inner sep=1pt,font=\fontsize{7}{8}\selectfont] at (0,-.20)
  {$v_{#1}(0)=#3\,\mathrm{L},\;s_{#1}(0)=\InitialSalt\,\mathrm{g}$};}
% One cylindrical tank: index, horizontal position, capacity, initial volume,
% initial concentration. The liquid height follows initial volume/capacity.
\newcommand{\ODETank}[5]{%
  \begin{scope}[shift={(#2,0)}]
    \pgfmathsetmacro{\ODElevel}{2.6*(#4)/(#3)}
    \path[fill=mtuGreenPale] (-.9,0) arc[start angle=180,end angle=360,x radius=.9,y radius=.17]
      -- (.9,\ODElevel) arc[start angle=0,end angle=180,x radius=.9,y radius=.17] -- cycle;
    \draw[mtuGreenLight,fill=mtuGreenLight!20] (0,\ODElevel) ellipse[x radius=.9,y radius=.17];
    \draw[thick] (-.9,2.6)--(-.9,0) arc[start angle=180,end angle=360,x radius=.9,y radius=.17] -- (.9,2.6);
    \draw[thick] (0,2.6) ellipse[x radius=.9,y radius=.17];
    \ODETankHeading{#1}{#3}
    \ODETankLabels{#1}{#3}{#4}{#5}
    \coordinate (tank#1wh) at (-.9,1.95);
    \coordinate (tank#1eh) at (.9,1.95);
    \coordinate (tank#1wl) at (-.9,.7);
    \coordinate (tank#1el) at (.9,.7);
  \end{scope}}

% Optional independent continuous salt rates in g/min. Zero disables each source.
% Add these rates to the salt-input vector, not to initial salt masses.
\newcommand{\ODETankSaltSource}[2]{%
  \ifdim#2pt>0pt\relax
    \begin{scope}[shift={(#1,0)},xshift=0.15cm,yshift=-0.12cm]
      % Belt sits below the capacity heading and delivers inward over the rim.
      \node[anchor=west,font=\small] at (1.12,3.20) {$#2\,\mathrm{g/min}$};
      \draw[thick,mtuGreen,rounded corners=2pt,fill=white]
        (0.48,2.80) rectangle (1.92,2.96);
      \draw[mtuGreen] (0.60,2.88) circle[radius=0.055];
      \draw[mtuGreen] (1.80,2.88) circle[radius=0.055];
      \draw[-{Stealth[length=1.5mm]},mtuGreen] (1.55,2.88)--(0.90,2.88);
      \foreach \sx/\sy in {1.65/3.00,1.43/3.00,1.20/3.00,0.98/3.00,0.75/3.00,0.50/2.75,0.48/2.64,0.53/2.53,0.48/2.42}{%
        \fill[red] (\sx,\sy) circle[radius=0.022];}
    \end{scope}
  \fi}
\newcommand{\ThreeTankDiagram}[1]{%
  \begingroup
  \pgfkeys{/ode tanks/.cd,
    capacity1=100,capacity2=120,capacity3=100,
    volume1=60,volume2=80,volume3=50,
    concentration1=2,concentration2=1,concentration3=0,
    inflow=2,inflowconcentration=3,flow12=5,flow21=3,flow23=4,flow32=2,outflow=2,
    saltRate1=0,saltRate2=0,saltRate3=0,#1}
  \begin{tikzpicture}[font=\small,
    pipe/.style={-{Stealth[length=2.5mm,width=1.7mm]},thick,mtuGreen},
    rate/.style={above,align=center,text=black,fill=white,inner sep=3pt}]
    \ODETank{1}{0}{\ODEcapone}{\ODEvolone}{\ODEcone}
    \ODETank{2}{4.6}{\ODEcaptwo}{\ODEvoltwo}{\ODEctwo}
    \ODETank{3}{9.2}{\ODEcapthree}{\ODEvolthree}{\ODEcthree}
    \ODETankSaltSource{0}{\ODEsaltone}
    \ODETankSaltSource{4.6}{\ODEsalttwo}
    \ODETankSaltSource{9.2}{\ODEsaltthree}
    \draw[pipe] (-3.7,1.95)--(tank1wh) node[midway,rate] {$\begin{aligned}q_{\mathrm{in}}&=\ODEqin\ \mathrm{L/min}\\c_{\mathrm{in}}&=\ODEcin\ \mathrm{g/L}\end{aligned}$};
    \draw[pipe] (tank1eh)--(tank2wh) node[midway,rate] {$q_{12}=\ODEqonetwo\ \mathrm{L/min}$};
    \draw[pipe] (tank2wl)--(tank1el) node[midway,rate] {$q_{21}=\ODEqtwoone\ \mathrm{L/min}$};
    \draw[pipe] (tank2eh)--(tank3wh) node[midway,rate] {$q_{23}=\ODEqtwothree\ \mathrm{L/min}$};
    \draw[pipe] (tank3wl)--(tank2el) node[midway,rate] {$q_{32}=\ODEqthreetwo\ \mathrm{L/min}$};
    \draw[pipe] (tank3el)--(12.9,.7) node[midway,rate] {$q_{\mathrm{out}}=\ODEqout\ \mathrm{L/min}$};
  \end{tikzpicture}%
  \endgroup}

\usepackage{bm}
\newcommand{\mat}[1]{\bm{#1}}
\newcommand{\SystemBlank}[2]{%
  \vcenter{\hbox{$\boxed{\rule{0pt}{#2}\makebox[#1]{}}$}}}
\newcommand{\VectorBlank}{\SystemBlank{3.2in}{0.80in}}
\newcommand{\MatrixBlank}{\SystemBlank{1.8in}{0.80in}}
\newcommand{\TrigMatrixBlank}{\SystemBlank{3.0in}{1.60in}}
\newcommand{\TimeBlank}[1][]{%
  \vcenter{\hbox{$\boxed{\parbox[c][0.25in][c]{0.45in}{\centering\ensuremath{#1}}}$}}}

\newcommand{\KeyBox}[3]{%
  \vcenter{\hbox{$\boxed{%
    \rule{0pt}{#2}\makebox[#1][c]{%
      \smash{\raisebox{\dimexpr#2/2\relax}{%
        \small\setlength{\arraycolsep}{3pt}\renewcommand{\arraystretch}{0.95}%
        $\vcenter{\hbox{$\displaystyle #3$}}$}}}}$}}}
\newcommand{\KeySide}[2]{\quad\vcenter{\hbox{\parbox{#1}{\footnotesize\setlength{\arraycolsep}{2pt}\ensuremath{\begin{gathered}#2\end{gathered}}}}}}
\begin{document}
\WorksheetHeading{Exam One Key}
\begin{enumerate}
\item[] \mbox{}
\newline
\resizebox{\linewidth}{!}{\ThreeTankDiagram{%
  capacity1=11,volume1=8,concentration1=1.5,capacity2=13,volume2=6,concentration2=0.5,capacity3=14,volume3=7,concentration3=1,inflow=4,inflowconcentration=7,flow12=7,flow21=3,flow23=6,flow32=2,outflow=4,saltRate1=5,saltRate2=6,saltRate3=8
}}
\newline
Tanks are well-mixed, $t$ in min, $\vect{v}(t)$ in L, $\vect{s}(t)$ in g, and conveyers add dry salt at the rate shown. 
\begin{enumerate}
\item
$\vect{s}'=\mat{A}\vect{s}+\vect{b}$,
$\vect{s}(0)=\begin{pmatrix}12\\3\\7\end{pmatrix}\,\mathrm{g}$\quad
$0<t<\TimeBlank[\infty]\quad\mathrm{min}$

\noindent $\vect{v}(t)=\KeyBox{1.6in}{0.80in}{\begin{pmatrix}8\\6\\7\end{pmatrix}}\quad\mathrm{L}\qquad
\vect{b}=\KeyBox{1.6in}{0.80in}{\begin{pmatrix}33\\6\\8\end{pmatrix}}\quad\mathrm{g/min}$
\item $\mat{A}=\KeyBox{3.0in}{1.60in}{\begin{pmatrix}-7/8&3/6&0\\7/8&-9/6&2/7\\0&6/6&-6/7\end{pmatrix}}\quad\mathrm{min}^{-1}\KeySide{2.5in}{\text{Brine and dry salt inputs (g/min)}\\ b_1=4(7)+5=33,\quad b_2=6,\quad b_3=8}$
\item A UC form is $\vect{s}_p(t)=\KeyBox{3.2in}{0.80in}{\vect{a}}$.\hfill{\footnotesize $\mat{A}\vect{a}=-\vect{b}$} 
Constants are determined by row reducing
$\mat{M}_{\text{\textsc{uc}}}=\KeyBox{1.8in}{0.80in}{\left[\begin{array}{rrr|r}-7/8&3/6&0&-33\\7/8&-9/6&2/7&-6\\0&6/6&-6/7&-8\end{array}\right]}\KeySide{3.4in}{\text{Net liquid flow by tank (L/min)}\\4+3-7=0,\quad7+2-3-6=0\\6-2-4=0}$

\end{enumerate}


\newpage
\item[] \mbox{}
\begin{enumerate}
\item \fbox{\normalsize\begin{tabular}{@{}l@{}}{\large$\rlap{\checkmark}\square$} Linear\\{\large$\square$} Separable\end{tabular}}\qquad $ty'-3y=t^4e^{2t}\qquad y(1)=1$
\par\noindent\makebox[2.7in][l]{Show a suitable intermediate step}\ $\KeyBox{3.2in}{0.80in}{(t^{-3}y)'=e^{2t}}$
\par\hfill{\footnotesize $\begin{gathered}y'-\frac{3}{t}y=t^3e^{2t},\quad \mu(t)=e^{-3\ln t}=t^{-3}\\t^{-3}y=\tfrac12 e^{2t}+C,\quad C=1-\tfrac12 e^2\end{gathered}$}\par
\vfill
\noindent\makebox[2.7in][l]{Show a fully integrated answer}\ $\KeyBox{3.2in}{0.80in}{t^{-3}y=1+\tfrac12(e^{2t}-e^2)}$\par
\par\medskip\hrule\medskip
\item \fbox{\normalsize\begin{tabular}{@{}l@{}}{\large$\square$} Linear\\{\large$\rlap{\checkmark}\square$} Separable\end{tabular}}\qquad $y'=3t^2\cos(t^3)(1+y)^2\qquad y(0)=1$
\par\noindent\makebox[2.7in][l]{Show a suitable intermediate step}\ $\KeyBox{3.2in}{0.80in}{\int\frac{dy}{(1+y)^2}=\int 3t^2\cos(t^3)\,dt}$
\par\hfill{\footnotesize $\begin{gathered}u=t^3,\quad du=3t^2\,dt\\-\frac{1}{1+y}=\sin(t^3)+C,\quad y(0)=1\ \Longrightarrow\ C=-\tfrac12\end{gathered}$}\par
\vfill
\noindent\makebox[2.7in][l]{Show a fully integrated answer}\ $\KeyBox{3.2in}{0.80in}{-\frac{1}{1+y}=\sin(t^3)-\tfrac12}$\par

\end{enumerate}
\newpage
\item[]
$\vect{y}'=\mat{A}\vect{y}+e^{3t}\begin{pmatrix}2\\-1\end{pmatrix}\quad
\mat{A}=\begin{pmatrix}1&-2\\8&1\end{pmatrix}\quad
\vect{y}(0)=\begin{pmatrix}1\\3\end{pmatrix}$
\par Eigenpairs $1+4i,\,\begin{pmatrix}1\\-2i\end{pmatrix}\quad
1-4i,\,\begin{pmatrix}1\\2i\end{pmatrix}$
\begin{enumerate}
\item $\vect{y}_c(t)=\KeyBox{3.2in}{0.80in}{\resizebox{3.05in}{!}{$\displaystyle c_1e^t\binom{\cos(4t)}{2\sin(4t)}+c_2e^t\binom{\sin(4t)}{-2\cos(4t)}\quad\text{or}\quad e^t\begin{pmatrix}\cos(4t)&\sin(4t)\\2\sin(4t)&-2\cos(4t)\end{pmatrix}\binom{c_1}{c_2}$}}\KeySide{2.45in}{\vect{v}=\binom{1}{0}+i\binom{0}{-2}\\e^{(1+4i)t}=e^t\cos(4t)+i e^t\sin(4t)\\\operatorname{re}\!\bigl(\vect{v}e^{(1+4i)t}\bigr)=e^t\binom{\cos(4t)}{2\sin(4t)}\\\operatorname{im}\!\bigl(\vect{v}e^{(1+4i)t}\bigr)=e^t\binom{\sin(4t)}{-2\cos(4t)}}$
\item A suitable UC form is
\newline
$\vect{y}_p(t)=\KeyBox{3.2in}{0.80in}{e^{3t}\vect{a}}\KeySide{2.45in}{\vect{y}_p'=3e^{3t}\vect{a}\\ \mat{A}\vect{y}_p+e^{3t}\binom{2}{-1}\\=e^{3t}\left(\mat{A}\vect{a}+\binom{2}{-1}\right)}$
\item To determine the UC coefficients you would row reduce
\newline
$\mat{M}_{\text{\textsc{uc}}}=\KeyBox{1.8in}{0.80in}{\left[\begin{array}{rr|r}2&2&2\\-8&2&-1\end{array}\right]}\KeySide{3in}{3\mat{I}-\mat{A}=\begin{pmatrix}3&0\\0&3\end{pmatrix}-\begin{pmatrix}1&-2\\8&1\end{pmatrix}\\=\begin{pmatrix}2&2\\-8&2\end{pmatrix}}$
\item To satisfy the IC you would row reduce
\par\noindent $\mat{M}_{\text{\textsc{ic}}}=\KeyBox{1.8in}{0.80in}{\left[\begin{array}{rr|r}1&0&1-a_1\\0&-2&3-a_2\end{array}\right]}\KeySide{3.2in}{\vect{y}_c(0)=\binom{c_1}{-2c_2},\quad \vect{y}_p(0)=\binom{a_1}{a_2}\\\vect{y}_c(0)+\vect{y}_p(0)=\binom{1}{3}}$\par
\end{enumerate}

\newpage
\item[]
$\vect{y}'=\mat{A}\vect{y}+\sin(3t)\begin{pmatrix}2\\1\end{pmatrix}\quad
\mat{A}=\begin{pmatrix}2&1\\0&-1\end{pmatrix}$
\par Eigenpairs $2,\,\begin{pmatrix}1\\0\end{pmatrix}\quad
-1,\,\begin{pmatrix}1\\-3\end{pmatrix}$
\begin{enumerate}
\item $\vect{y}_c(t)=\KeyBox{3.2in}{0.80in}{\resizebox{3.05in}{!}{$\displaystyle c_1e^{2t}\begin{pmatrix}1\\0\end{pmatrix}+c_2e^{-t}\begin{pmatrix}1\\-3\end{pmatrix}\quad\text{or}\quad \begin{pmatrix}e^{2t}&e^{-t}\\0&-3e^{-t}\end{pmatrix}\binom{c_1}{c_2}$}}$
\item A suitable UC form is
\newline
$\vect{y}_p(t)=\KeyBox{3.2in}{0.80in}{\vect{a}\cos(3t)+\vect{b}\sin(3t)}\KeySide{2.45in}{\vect{y}_p'=3\vect{b}\cos(3t)-3\vect{a}\sin(3t)\\\mat{A}\vect{y}_p+\binom{2}{1}\sin(3t)\\=(\mat{A}\vect{a})\cos(3t)\\+\left(\mat{A}\vect{b}+\binom{2}{1}\right)\sin(3t)}$
\item To determine the UC coefficients you would row reduce
\newline
$\mat{M}_{\text{\textsc{uc}}}=\KeyBox{2.7in}{1.20in}{\left[\begin{array}{c@{\qquad}c|c}-\mat{A}&3\mat{I}&\begin{pmatrix}0\\0\end{pmatrix}\\[10pt]-3\mat{I}&-\mat{A}&\begin{pmatrix}2\\1\end{pmatrix}\end{array}\right]}\KeySide{3in}{-\mat{A}\vect{a}+3\mat{I}\vect{b}=\binom{0}{0}\quad(\cos)\\-3\mat{I}\vect{a}-\mat{A}\vect{b}=\begin{pmatrix}2\\1\end{pmatrix}\quad(\sin)\\-\mat{A}=\begin{pmatrix}-2&-1\\0&1\end{pmatrix}\\\text{Columns: }a_1,a_2,b_1,b_2}$
\end{enumerate}
\end{enumerate}
\end{document}